\section{Related Work}
\label{sec:related}

\para{Model collapse under recursive training.}
\citet{shumailov2023curse,shumailov2024collapse} show that models trained repeatedly on the outputs of earlier model generations forget the tails of the original distribution.
\citet{alemohammad2024mad} report a similar loss of quality or diversity in self-consuming image generators when too little fresh real data enters each generation.
\citet{dohmatob2024tale} analyze collapse as a change in scaling laws, and \citet{guo2023curious} document declining linguistic diversity when language models train on synthetic text.
\citet{gerstgrasser2024inevitable} and \citet{bertrand2024stability} find that accumulating or retaining real data alongside synthetic data limits collapse.
Unlike this line of work, we study a single round of generation by a stronger teacher with a fixed real set, and we vary the mixing ratio rather than the number of generations.
We ask whether the diversity loss that drives recursive collapse is already visible, and already costly, after one round.

\para{Synthetic data for augmentation.}
In vision, \citet{he2023synthetic} and \citet{azizi2023synthetic} examine when images from generative models help recognition, and \citet{fan2024scaling} study how the benefit scales with the number of synthetic images.
In language, \citet{ye2022zerogen} train small task models entirely on generated data, \citet{yu2023attributed} show that attribute-controlled prompts make generated data more diverse and less biased, and \citet{li2023synthetic} find that the usefulness of generated data for text classification depends on the task.
Each of these studies covers one task family with its own generator, student, and metric.
Building on their observation that more synthetic data gives diminishing returns, we hold the teacher, student, real-set size, and ratio grid fixed across three task families, and we add a real-data control at matched size so that a plateau can be attributed to the data source and not to the capacity of the student.
The setting is also close to knowledge distillation~\citep{hinton2015distilling}, with generated examples carrying the teacher's knowledge instead of soft labels.

\para{Measuring diversity.}
Lexical diversity of generated text is commonly summarized by distinct-$n$~\citep{li2016diversity}.
The Vendi score~\citep{friedman2023vendi} gives an effective number of distinct items from the eigenvalues of a similarity matrix, and Fr\'echet distances between embedding distributions~\citep{heusel2017gans} compare a generated set to a reference set.
We use all three and find that they disagree in an informative way: synthetic-trained students keep their distinct-2 score while their stories become more alike in embedding space.
